% RAG-Verify -- IEEE conference draft (CONECCT 2027 target).
% DRAFT STATUS (28 Oct 2026): every number in this file comes from results/ (see code/aggregate.py,
% code/abstention_breakdown.py, results/tuning/selected.json). Items still running are marked \pending{} / \todo{}.
% Built on the official IEEE conference template (version 6/27/2024): IEEEtran.cls V1.8b and IEEEtran.bst V1.12,
% both copied into this folder from the IEEE template packages.
\documentclass[conference]{IEEEtran}
\IEEEoverridecommandlockouts
% The preceding line is only needed to identify funding in the first footnote. If that is unneeded, please comment it out.
\usepackage{cite}
\usepackage{amsmath,amssymb,amsfonts}
\usepackage{algorithmic}
\usepackage{graphicx}
\usepackage{textcomp}
\usepackage{xcolor}
% Additions to the template package list: a floating algorithm box, table rules, and URL line breaking.
\usepackage{algorithm}
\usepackage{booktabs}
\usepackage{url}
\def\BibTeX{{\rm B\kern-.05em{\sc i\kern-.025em b}\kern-.08em
    T\kern-.1667em\lower.7ex\hbox{E}\kern-.125emX}}

% \input inside a tabular breaks \multicolumn; the TeX primitive input does not.
\makeatletter\newcommand{\inputrows}{\@@input}\makeatother
\newcommand{\pending}{\textcolor{red}{\textsc{pend.}}}
\newcommand{\todo}[1]{\textcolor{red}{[PENDING: #1]}}

\begin{document}

\title{When Does Claim-Level Verification Help Small Retrieval-Augmented LLMs? An Empirical Study with RAG-Verify}

% Author names, order, affiliation and the corresponding author's e-mail as in the authors' Information Fusion
% preprint (28 Oct 2026); ORCID iDs supplied by the first author on 28 Oct 2026 (check digits verified).
\author{\IEEEauthorblockN{1\textsuperscript{st} Shivam Sengar}
\IEEEauthorblockA{\textit{Dept. of Computer Science and Engineering} \\
\textit{Defence Institute of Advanced Technology (DIAT)}\\
Pune, Maharashtra, India \\
shivamswngar4321@gmail.com \\
ORCID: 0009-0007-4037-9595}
\and
\IEEEauthorblockN{2\textsuperscript{nd} Bansidhar Joshi}
\IEEEauthorblockA{\textit{Dept. of Computer Science and Engineering} \\
\textit{Defence Institute of Advanced Technology (DIAT)}\\
Pune, Maharashtra, India \\
bansidhar11@gmail.com \\
ORCID: 0000-0001-7141-7258}
\and
\IEEEauthorblockN{3\textsuperscript{rd} Yash Prabhakar Sonawane}
\IEEEauthorblockA{\textit{Dept. of Computer Science and Engineering} \\
\textit{Defence Institute of Advanced Technology (DIAT)}\\
Pune, Maharashtra, India \\
sonawaney789@gmail.com \\
ORCID: 0009-0000-1556-4616}
}

\maketitle

\begin{abstract}
Retrieval-augmented generation (RAG) reduces hallucination, yet small open models still give confident wrong answers when the retrieved evidence is incomplete. A popular remedy is a post-hoc check: split the answer into atomic claims, test each claim against the evidence, repair what fails, and abstain when the answer cannot be supported. We ask whether this machinery earns its cost when the whole system must run on one 8\,GB consumer GPU. We build RAG-Verify, which combines a question-aware LLM check of the answer claim with natural language inference (NLI) on the remaining claims, followed by targeted revision and an abstention gate. We compare it with matched baselines on three disjoint 300-question test blocks of HotpotQA and TriviaQA using 4-bit Qwen2.5-7B and Llama-3.1-8B, with every design choice fixed beforehand on a separate development block. With Qwen2.5-7B, RAG-Verify cuts the hallucination rate of standard RAG from 50.6\% to 21.0\% on HotpotQA and from 23.8\% to 7.9\% on TriviaQA, and significantly raises effective reliability on both. However, a simpler baseline that asks the model to re-read the evidence and revise its answer once, with no verifier, is significantly better than RAG-Verify with Qwen and statistically tied with it with Llama, using fewer model calls. Trace analysis shows that the verifier flags 39\% of initially correct HotpotQA answers, and most flagged drafts end in abstention. At this scale, the benefit attributed to claim-level verification comes mainly from giving the model a second, evidence-focused pass.
\end{abstract}

\begin{IEEEkeywords}
large language models, hallucination, retrieval-augmented generation, fact verification, natural language inference, selective question answering
\end{IEEEkeywords}

% =====================================================================
\section{Introduction}
Large language models (LLMs) produce fluent text that is not always true, a failure usually called hallucination~\cite{huang2023survey,li2025capability}. Grounding the model in retrieved passages~\cite{lewis2020rag,gao2023ragsurvey} reduces the problem but does not remove it: the retriever can miss the decisive passage, and the model may ignore or misread evidence that is present~\cite{liu2023lost,niu2023ragtruth}. A growing family of methods therefore inspects the output after generation. Some plan and answer verification questions~\cite{dhuliawala2023cove}, some revise text using retrieved evidence~\cite{gao2022rarr}, some decompose an answer into atomic facts and test each one~\cite{min2023factscore,chern2023factool,wei2024safe,afev2025,vladika2025medical}, and others train the generator to critique its own retrievals~\cite{asai2023selfrag}.

Most of these studies use large proprietary models as the generator, the verifier, or both. Many practitioners, however, must run everything locally: on-premise deployments, privacy-sensitive settings, and teaching or research labs with a single consumer GPU. In that regime every extra model call costs latency, and the components that make verification work, namely claim decomposition, entailment checking, and faithful revision, are themselves carried out by small models that make mistakes. It is not obvious that stacking these imperfect steps helps.

This paper reports a controlled study of that question. We implement \emph{RAG-Verify}, a claim-level verify-and-revise pipeline built only from open models that fit in 8\,GB of GPU memory, and we compare it against baselines that share its retriever, prompts, and backbone but drop one ingredient at a time. Our contributions are:
\begin{enumerate}
\item A fully local claim-level verification pipeline with a \emph{hybrid} verifier: a question-aware LLM check for the claim that carries the answer and an NLI model for the supporting claims. We show on development data why plain NLI is not sufficient for the answer claim.
\item A pre-registered evaluation protocol: design, thresholds, and abstention rule are chosen on a held-out development block with a rule fixed in advance, and all systems are then run once on three disjoint test blocks per dataset, with paired bootstrap tests.
\item A negative but practically useful finding: an unconditional self-revision step with no verifier is significantly better than full claim-level verification on both datasets with Qwen2.5-7B and statistically tied with it with Llama-3.1-8B, with no more than 60\% of its model calls. A trace-level analysis locates the cause in verifier false alarms that convert correct answers into abstentions.
\end{enumerate}

% =====================================================================
\section{Related Work}
\textbf{Retrieval and adaptive retrieval.} RAG conditions generation on passages from a dense or sparse retriever~\cite{lewis2020rag,gao2023ragsurvey}. FLARE~\cite{jiang2023flare} decides when to retrieve during generation, CRAG~\cite{yan2024crag} grades retrieved documents and falls back to other sources when they look unreliable, and Self-RAG~\cite{asai2023selfrag} trains the generator to emit reflection tokens that score relevance, support, and usefulness of each passage.

\textbf{Post-hoc verification and revision.} Chain-of-Verification (CoVe)~\cite{dhuliawala2023cove} drafts an answer, plans verification questions, answers them independently, and writes a final response. RARR~\cite{gao2022rarr} retrieves evidence for an existing output and edits unsupported spans. Claim-level pipelines such as FActScore~\cite{min2023factscore}, FacTool~\cite{chern2023factool}, and SAFE~\cite{wei2024safe} split text into atomic facts and check each one against a knowledge source; more recent work refines the decomposition step~\cite{afev2025}, applies atomic checking to medical RAG~\cite{vladika2025medical}, separates retrieval, verification and calibration into modules~\cite{r2vc2026}, or verifies entities against evidence~\cite{jia2026entity}. Lightweight checkers such as MiniCheck~\cite{tang2024minicheck} and NLI models trained on many entailment datasets~\cite{laurer2023nli,he2021debertav3} make claim checking affordable.

\textbf{Detection without retrieval and abstention.} SelfCheckGPT~\cite{manakul2023selfcheckgpt} and semantic entropy~\cite{farquhar2024semantic} detect hallucinations from disagreement between samples. Benchmarks such as HaluEval~\cite{li2023halueval}, TruthfulQA~\cite{lin2021truthfulqa}, RAGTruth~\cite{niu2023ragtruth} and FaithJudge~\cite{tamber2025faithjudge} measure the problem, and recent work questions how detectors are evaluated~\cite{chen2026rethinking}. Declining to answer is a legitimate outcome when the model is unsure~\cite{wen2024abstention}; we adopt effective reliability~\cite{whitehead2022reliable}, which rewards correct answers, penalises wrong ones, and gives zero to abstentions.

Our study differs from the above in scope rather than mechanism: we hold the compute budget to one consumer GPU, include the simple baselines that a deployer would try first, and report what each added component contributes.

% =====================================================================
\section{Method}
\subsection{Pipeline}
Given a question $q$, a retriever returns the top-$k$ passages $P=\{p_1,\dots,p_k\}$. The generator $G$ produces a short response $r_0$ in a fixed format (answer line plus one-sentence explanation). RAG-Verify then runs the steps in Algorithm~\ref{alg:rv}.

\begin{algorithm}[t]
\caption{RAG-Verify (one question)}\label{alg:rv}
\begin{algorithmic}[1]
\REQUIRE question $q$, retriever $R$, generator $G$, NLI model $N$, threshold $\tau$, ratio $\rho$
\STATE $P \gets R(q,k)$; \quad $r_0 \gets G(\textsc{Answer}(q,P))$
\STATE $C \gets \textsc{Decompose}_G(q,r_0)$ \COMMENT{answer claim first}
\STATE $\ell \gets \textsc{Verify}(q,C,P)$ \COMMENT{hybrid, Eq.~(\ref{eq:label})}
\STATE $r_1 \gets r_0$
\IF{some $\ell_i \neq \textsc{supported}$}
  \STATE $r_1 \gets G(\textsc{Revise}(q,P,r_0,\{(c_i,\ell_i)\}))$
  \IF{$r_1$ is a refusal (``I don't know'')}
    \RETURN \textsc{abstain}
  \ENDIF
  \STATE $C \gets \textsc{Decompose}_G(q,r_1)$; \quad $\ell \gets \textsc{Verify}(q,C,P)$
\ENDIF
\IF{$\ell_1 \neq \textsc{supported}$ \textbf{or} $s(\ell) < \rho$}
  \RETURN \textsc{abstain}
\ENDIF
\RETURN the answer in $r_1$
\end{algorithmic}
\end{algorithm}

\textbf{Decomposition.} $G$ is prompted with two worked examples to rewrite the response as at most twelve self-contained atomic claims, the first of which must state the answer. Because small models do not always respect that ordering, we move the claim that contains the answer string to the front, or insert the template claim \emph{``The answer to the question `$q$' is $a$''} when no claim contains it.

\textbf{Hybrid verification.} For a supporting claim $c_i$ ($i>1$) the NLI model returns entailment and contradiction probabilities for each premise; premises are the $k$ passages plus the concatenation of the top three, which allows two-hop facts to be entailed. We take the maximum over premises, $e_i$ and $x_i$, and assign
\begin{equation}
\ell_i=\begin{cases}\textsc{supported} & e_i\ge\tau\\ \textsc{contradicted} & e_i<\tau,\ x_i\ge\tau\\ \textsc{unsupported} & \text{otherwise.}\end{cases}\label{eq:label}
\end{equation}
For the answer claim $c_1$ we do not use NLI. On the development block we observed that NLI entailed claims such as ``the answer is $X$'' whenever $X$ appeared in a passage, even when $X$ answered a different question; it caught only 3 of 17 wrong answers in a manual check, against 7 of 17 for an LLM judge. The answer claim is therefore judged by $G$ itself with a prompt that asks whether the proposed answer is correct \emph{for this question} given the passages, returning one of the three labels.

\textbf{Revision and abstention.} Flagged claims and their labels are shown to $G$ together with the passages and the original response; $G$ may correct the answer or reply that it does not know. After re-verification, the system abstains if the answer claim is not supported or if the supported fraction
\begin{equation}
s(\ell)=\frac{1}{|C|}\sum_{i=1}^{|C|}\mathbf{1}[\ell_i=\textsc{supported}]
\end{equation}
falls below $\rho$.

\subsection{Design selection on development data}
Before any test run we fixed the selection rule: maximise mean effective reliability over the two development blocks (100 questions each), breaking ties within 0.01 by higher coverage, then fewer model calls, then $\tau$ closest to 0.5. We compared three verifiers (NLI only, hybrid, LLM judge for every claim), $\tau\in\{0.3,0.5,0.7,0.9\}$, and $\rho\in[0,1]$ (simulated offline from saved scores). The rule selected the hybrid verifier with $\tau=0.3$ and $\rho=0$, i.e., abstention is driven by the answer claim alone; mean development reliability was 0.200. This configuration was frozen and is the one reported as RAG-Verify.

% =====================================================================
\section{Experimental Setup}
\textbf{Data.} We use HotpotQA (distractor setting, validation split)~\cite{yang2018hotpotqa} and the Wikipedia subset of TriviaQA (validation split)~\cite{joshi2017triviaqa}. A fixed permutation (seed 2026) of each split gives three disjoint 300-question test blocks and one 100-question development block; we report test results only. Retrieval corpora are built from the provided contexts: 66{,}635 unique HotpotQA paragraphs from all 7{,}405 validation questions, and 135{,}517 100-word passages from 3{,}120 TriviaQA evidence pages (the 1{,}000 sampled questions plus 1{,}100 distractor questions). These are closed corpora, not full Wikipedia, so absolute scores are not comparable to open-domain results. To test the verifiers in isolation we also use HaluEval-QA~\cite{li2023halueval}: 1{,}000 questions sampled with the same seed, each with its gold knowledge snippet and one correct and one hallucinated answer (2{,}000 items).

\textbf{Models.} Generators are Qwen2.5-7B-Instruct~\cite{qwen2024qwen25} and Llama-3.1-8B-Instruct~\cite{dubey2024llama3}, both 4-bit (Q4\_K\_M) served by Ollama with temperature 0, a fixed seed, and an 8{,}192-token context. The retriever is bge-small-en-v1.5~\cite{xiao2023cpack} with exact inner-product search, $k=5$. The NLI model is DeBERTa-v3-large~\cite{he2021debertav3} fine-tuned on MNLI, FEVER~\cite{thorne2018fever}, ANLI, LingNLI and WANLI~\cite{laurer2023nli}. Self-RAG uses the authors' 7B checkpoint~\cite{asai2023selfrag} quantised to 4-bit NF4. Hardware is a laptop with an RTX~4060 (8\,GB).

\textbf{Systems.} All retrieval-based systems share the same retriever, passages, and answer format. \emph{Closed-book} answers without passages. \emph{Standard RAG} answers from the passages. \emph{RAG + abstain prompt} adds an instruction to reply ``I don't know'' when the passages are insufficient. \emph{RAG + self-revision} shows the model its own first response and the passages and asks it to correct any error or abstain, with no verifier signal; it is RAG-Verify with the verifier removed. \emph{CoVe}~\cite{dhuliawala2023cove} uses the factored variant with up to five verification questions, either closed-book or with retrieval for every verification question. \emph{Self-RAG}~\cite{asai2023selfrag} follows the authors' short-form inference (adaptive retrieval threshold 0.2, equal weights for the critique scores).

\textbf{Metrics.} An answer is correct if, after SQuAD normalisation~\cite{rajpurkar2016squad}, it equals or contains a gold alias. Let $A$ be accuracy over all questions and $H$ the fraction of questions answered incorrectly (hallucination rate); abstentions are neither. Effective reliability~\cite{whitehead2022reliable} with unit cost is
\begin{equation}
\Phi_1 = A - H = \frac{1}{n}\sum_{j=1}^{n}\begin{cases}+1 & \text{correct}\\ -1 & \text{answered, wrong}\\ 0 & \text{abstained.}\end{cases}
\end{equation}
We report mean $\pm$ standard deviation over the three test blocks, and paired bootstrap 95\% confidence intervals (10{,}000 resamples) over the 900 pooled questions. On the first test block we also measure FActScore-style claim precision of the answered outputs: a fixed extractor (Llama-3.1-8B, the same for every system) splits each output into atomic claims, and MiniCheck-Flan-T5-Large~\cite{tang2024minicheck} judges each claim against gold evidence (the HotpotQA supporting paragraphs; for TriviaQA the five chunks of the gold Wikipedia pages closest to the claim), never against the pipeline's own retrieval. Latency is the median wall-clock time per question on the first 50 development questions, with one system on the GPU at a time.

% =====================================================================
\section{Results}
\begin{table*}[t]
\centering
\caption{Test results (\%, mean $\pm$ std over three disjoint 300-question blocks). Acc: accuracy; Abst: abstention rate; Hall: answered and wrong; $\Phi_1$: effective reliability (higher is better).}
\label{tab:main}
\setlength{\tabcolsep}{4pt}
\begin{tabular}{lcccccccc}
\toprule
 & \multicolumn{4}{c}{HotpotQA} & \multicolumn{4}{c}{TriviaQA} \\
\cmidrule(lr){2-5}\cmidrule(lr){6-9}
System & Acc$\uparrow$ & Abst & Hall$\downarrow$ & $\Phi_1\uparrow$ & Acc$\uparrow$ & Abst & Hall$\downarrow$ & $\Phi_1\uparrow$ \\
\midrule
\inputrows tables/main_rows.tex 
\bottomrule
\end{tabular}
\end{table*}

\begin{table}[t]
\centering
\caption{Paired bootstrap differences in $\Phi_1$ (percentage points, 95\% CI, 900 questions per dataset). $^{*}$: CI excludes zero.}
\label{tab:tests}
\setlength{\tabcolsep}{2pt}
\resizebox{\columnwidth}{!}{%
\begin{tabular}{lcc}
\toprule
Comparison & HotpotQA & TriviaQA \\
\midrule
\inputrows tables/tests_rows.tex 
\bottomrule
\end{tabular}}
\end{table}

Table~\ref{tab:main} lists all systems and Table~\ref{tab:tests} the paired comparisons.

\textbf{Retrieval matters most.} With Qwen2.5-7B, retrieval raises accuracy from 21.9\% to 47.4\% on HotpotQA and from 57.2\% to 75.6\% on TriviaQA. Closed-book CoVe does not improve on the closed-book model (21.8\% and 60.7\%), and CoVe with retrieval is close to standard RAG ($\Phi_1$ $-$2.0 vs.\ $-$3.1 on HotpotQA, 54.8 vs.\ 51.8 on TriviaQA) while making about 7.7 model calls per question. Verification questions answered by the same small model rarely overturn its first answer.

\textbf{RAG-Verify trades coverage for precision.} RAG-Verify lowers the hallucination rate of standard RAG from 50.6\% to 21.0\% on HotpotQA and from 23.8\% to 7.9\% on TriviaQA, the lowest of all systems on both datasets. It abstains on 49.2\% and 27.4\% of questions, so accuracy falls to 29.8\% and 64.7\%. The net effect on $\Phi_1$ is positive and significant: $+$11.9 points on HotpotQA and $+$5.0 on TriviaQA. When it does answer, RAG-Verify is right 58.5\% (HotpotQA) and 89.1\% (TriviaQA) of the time, against 48.4\% and 76.1\% for standard RAG.

\textbf{Simple baselines are as good or better.} A one-line abstention instruction already reaches $\Phi_1$ of 12.0 and 56.2; the difference from RAG-Verify is not significant on either dataset. Unconditional self-revision does better still: 15.9 on HotpotQA and 59.9 on TriviaQA, significantly above both RAG-Verify ($-$7.1 and $-$3.1 points for RAG-Verify) and the abstention prompt ($+$3.9 and $+$3.7). Self-revision makes exactly two model calls per question, whereas RAG-Verify makes 3.7 (HotpotQA) and 3.4 (TriviaQA) calls on average plus NLI inference.

\textbf{Backbone.} With Llama-3.1-8B the closed-book model is much stronger on TriviaQA (74.9\% vs.\ 57.2\%), so retrieval adds less (77.4\%), and the abstention prompt is followed less often (13.9\% and 9.9\% abstention vs.\ 34.0\% and 22.9\% for Qwen). RAG-Verify again beats standard RAG significantly ($\Phi_1$ $+$7.4 and $+$4.3 points) and again has the lowest hallucination rate among Llama systems (31.4\% and 11.6\%). Unlike with Qwen, however, self-revision is not better than RAG-Verify: the two are statistically tied on both datasets (8.1 vs.\ 6.4 on HotpotQA, 61.8 vs.\ 61.4 on TriviaQA), as are RAG-Verify and the abstention prompt. Llama abstains less after a flag, so the verifier costs fewer correct answers, but RAG-Verify now needs 5.2 and 4.9 model calls per question against two for self-revision. Self-RAG, which answers almost every question (0.8\% and 0.6\% abstention), has the highest hallucination rate of the retrieval-based systems (59.4\% and 25.7\%) and $\Phi_1$ of $-$19.7 and 48.1; this comparison is only indicative, since its backbone and quantisation differ from ours.

% =====================================================================
\section{Analysis}
\begin{table}[t]
\centering
\caption{Where RAG-Verify's decisions come from (Qwen2.5-7B, 900 test questions per dataset), computed from saved traces.}
\label{tab:trace}
\begin{tabular}{lrr}
\toprule
 & HotpotQA & TriviaQA \\
\midrule
First draft correct & 425 & 682 \\
First draft flagged by verifier & 466 & 282 \\
\quad of which first draft was correct & 164 & 134 \\
Abstained: model refused in first draft & 19 & 7 \\
Abstained: model refused after revision & 380 & 233 \\
Abstained: final gate (answer claim) & 44 & 7 \\
\quad of which withheld answer was correct & 24 & 3 \\
Answered & 457 & 653 \\
\bottomrule
\end{tabular}
\end{table}

Table~\ref{tab:trace} breaks RAG-Verify's behaviour down using the per-question traces. Two observations explain the gap to self-revision.

\emph{The verifier raises many false alarms.} On HotpotQA it flags 466 of 900 first drafts, and 164 of those drafts were already correct; that is, 38.6\% of correct first answers are sent back for revision (19.6\% on TriviaQA). The error analysis below shows that these flags come mostly from the question-aware check of the answer claim, not from NLI.

\emph{A flag pushes the model towards refusal.} Most abstentions (380 of 443 on HotpotQA, 233 of 247 on TriviaQA) come from the model itself choosing ``I don't know'' after being told that some of its claims are unsupported; equivalently, 82\% of flagged HotpotQA drafts (380 of 466) and 83\% of flagged TriviaQA drafts (233 of 282) end in a refusal; the final gate adds only a few, and more than half of those it adds on HotpotQA (24 of 44) withhold a correct answer. Unconditional self-revision asks the model to re-check every answer without signalling that anything is wrong, and the model keeps more of its correct answers while still dropping many wrong ones.

The development-block comparison of verifier types pointed the same way: the hybrid verifier was selected by the pre-declared rule, but it was not significantly different from self-revision on development data either ($-$0.025, 95\% CI [$-$0.09, 0.04]). We froze the design before testing and report the outcome as observed.

\begin{table}[t]
\centering
\caption{Ablations on the first test block only (300 questions per dataset, Qwen2.5-7B). $\Phi_1$ in \%; $\Delta$: paired difference from the reference (full RAG-Verify; standard RAG for the last group) with 95\% bootstrap CI. $^{*}$: CI excludes zero. The $\tau$ and $k$ rows are sensitivity checks only.}
\label{tab:abl}
\setlength{\tabcolsep}{2pt}
\resizebox{\columnwidth}{!}{%
\begin{tabular}{lcccc}
\toprule
 & \multicolumn{2}{c}{HotpotQA} & \multicolumn{2}{c}{TriviaQA} \\
\cmidrule(lr){2-3}\cmidrule(lr){4-5}
Variant & $\Phi_1$ & $\Delta$ [95\% CI] & $\Phi_1$ & $\Delta$ [95\% CI] \\
\midrule
\inputrows tables/ablation_rows.tex
\bottomrule
\end{tabular}}
\end{table}

\textbf{Ablations.} Table~\ref{tab:abl} removes or varies one component at a time. Retrieval and the abstention gate carry the effect. Without retrieval, where claims are checked by agreement between five sampled answers as in SelfCheckGPT~\cite{manakul2023selfcheckgpt}, $\Phi_1$ falls by 25.3 and 19.7 points; without the gate the hallucination rate returns to that of standard RAG (50.0\% vs.\ 47.3\% on HotpotQA, 25.7\% for both on TriviaQA). Removing claim decomposition or the revision step changes $\Phi_1$ by at most one point, and $\tau$ matters only at 0.9 on TriviaQA. This is expected from the selected $\rho=0$: abstention depends only on the answer claim, which the LLM judges, so the NLI-checked supporting claims can trigger a revision but never block an answer. Using NLI or the LLM judge for every claim is no better than the hybrid. Retrieving one passage hurts both RAG-Verify and standard RAG; three or ten passages make no significant difference to RAG-Verify. Removing the verifier altogether is the only change with a significant gain ($+$10.3 points on HotpotQA).

\textbf{Error analysis.} We assigned every failure of RAG-Verify to a cause with rules over the traces. A \emph{lost answer} is an abstention after a correct first draft: there are 156 on HotpotQA and 100 on TriviaQA, and the answer-claim judge rejected the correct answer in 149 and 87 of them. Of the 189 and 71 wrong answers, 60 and 22 had no gold answer string in the retrieved passages, and in 82 and 37 the judge accepted a wrong answer although the gold answer had been retrieved. We then read a random sample of 50 failures (25 per dataset; one reading, not cross-checked by a second annotator). Of the 25 lost answers, 9 were rejected although the passages stated the answer and 5 were comparison or yes/no questions that need two passages; in 7 the draft was right but drawn from the model's memory rather than the passages, where abstaining is defensible, and 3 were artefacts of noisy TriviaQA aliases. Of the 25 wrong answers, 9 named a related entity instead of the attribute asked for (e.g., a singer's name for a question about the singer's birth date) and were accepted by the judge, 6 were reasoning or reading errors, 4 followed a retrieval miss, and 6 were acceptable answers missed by string matching. The judge thus errs in both directions: it checks whether an answer is mentioned and plausible, not whether it fills the slot the question asks for.

\textbf{Verifiers in isolation.} On HaluEval-QA, with the gold knowledge as evidence and each answer phrased as the same question-grounded claim, NLI ranks hallucinated above correct answers with AUROC 0.876 (95\% CI 0.860--0.890). At $\tau=0.3$, however, it flags only 607 of the 1{,}000 hallucinated answers, with 40 false alarms on correct ones. The LLM judge is more balanced (accuracy 0.810, F1 0.805 vs.\ 0.737 for NLI) but flags 164 of the 1{,}000 correct answers even though the evidence is clean and short. A false-alarm rate of this size on correct answers is in line with the 19.6\% of correct TriviaQA drafts that the hybrid verifier flags in the main experiment.

\textbf{Claim precision and cost.} On the first test block, 72.5\% (HotpotQA) and 71.2\% (TriviaQA) of the claims in RAG-Verify's answered outputs are supported by gold evidence, against 62.0\% and 61.5\% for standard RAG, 68.5\% and 67.7\% for self-revision, 69.7\% and 69.4\% for the abstention prompt, 58.4\% and 59.5\% for CoVe with retrieval, and 27.2\% and 45.2\% closed-book. Precision is computed only over answered questions, so systems that abstain more are measured on easier subsets; these are single-block estimates without intervals. The median time per question is 3.7\,s and 5.2\,s for RAG-Verify, 2.0\,s for self-revision, 1.1\,s for standard RAG, and 6.9\,s and 7.6\,s for CoVe with retrieval.

% =====================================================================
\section{Discussion and Limitations}
Our results do not show that claim-level verification is useless; it produced the lowest hallucination rates in the study, which may be the right operating point when a wrong answer is far more costly than a refusal. They do show that, with 7--8B generators on a consumer GPU, most of its benefit under a symmetric cost can be obtained by a single, unguided revision pass, and that a noisy verifier converts many correct answers into refusals. Better verifiers, verifier confidence rather than hard labels, or asymmetric costs could change the picture.

The study has clear limits. Test sets are 900 questions per dataset drawn from two short-answer QA benchmarks, and the retrieval corpora are closed. Answer correctness uses string matching, which can miss valid paraphrases. Generation runs through a quantised inference server that is not bit-reproducible: in our checks, repeated runs gave identical final answers but occasionally different explanation text, and the server was updated from version 0.40.11 to 0.40.12 during the experiments (final answers on a 60-question probe were unchanged). Self-RAG uses a 4-bit checkpoint and our retriever, which differs from the original setup. Claim precision relies on an automatic extractor and MiniCheck, which we did not validate against human judgements on our data.

% =====================================================================
\section{Conclusion}
We built a fully local claim-level verify-and-revise pipeline and tested it against matched baselines under a pre-registered protocol. It reliably reduces hallucinations relative to standard RAG, but an unconditional self-revision step is as good in effective reliability (Llama-3.1-8B) or significantly better (Qwen2.5-7B), and cheaper. For small open models, the second look at the evidence appears to matter more than the fine-grained verdicts that are supposed to guide it. Future work should test whether calibrated verifier scores, rather than discrete flags, can recover the correct answers that the current design discards.

\section*{Reproducibility}
Code, configurations, prompts, per-question outputs, and logs are in the project repository at \url{https://github.com/shivamsengar135/RAG-Verify}; every table is regenerated by \texttt{code/aggregate.py} from the saved outputs.

\bibliographystyle{IEEEtran}
\bibliography{refs}

\end{document}
